% =====================================================================
% POSITIONING
%
% Scope note (mirrors the English edition): this section positions the
% study against the two papers we hold and have read in full. It does not
% survey the wider literature, because a survey assembled from titles and
% abstracts would be a list of things we had not read. Tools and
% taxonomies the study uses -- Garak, the OWASP LLM Top 10 -- are named
% where they are used and addressed in a footnote, not cited as sources
% of claims.
% =====================================================================
\section{جایگاه‌یابی}
\label{sec:related}

\subsection{\en{Prompt Injection} غیرمستقیم}
\label{sec:rel:injection}

\en{Greshake} و همکاران~\cite{greshake2023} پیشینهٔ مستقیمِ نیمهٔ تزریقِ
این مطالعه‌اند. حرکت آن‌ها این است که بپرسند وقتی \emph{کاربر} نیست که
\en{prompt} می‌دهد چه اتفاقی می‌افتد: در یک برنامهٔ یکپارچه با \en{LLM}،
مهاجم می‌تواند دستورها را در داده‌ای بگذارد که برنامه در زمان استنتاج
بازیابی خواهد کرد و بدین ترتیب مدل را کنترل کند بی‌آنکه هرگز رابطی به آن
داشته باشد. نتیجه‌ای که می‌گیرند، یعنی اینکه پردازش محتوای بازیابی‌شدهٔ
غیرقابل‌اعتماد قیاس‌پذیر با اجرای کد دلخواه است، چون مرزِ میان داده و
دستور در چنین برنامه‌ای هیچ‌جای سامانه کشیده
نشده~\cite[\en{Sec.~2}]{greshake2023}، همان مقدمه‌ای است که این مقاله از
آن آغاز می‌کند و سپس روی یک برنامه اندازه‌اش می‌گیرد.

سه بخش از کار آن‌ها به‌طور مشخص به کار می‌رود، نه در گذر. نخست،
\emph{روش‌های تزریقِ} آن‌ها~\cite[\en{Sec.~3.1}]{greshake2023} تحویلِ
منفعل از راه \en{retrieval}، تحویلِ فعال مانند \en{email}، تزریقِ
کاربر-محور که در آن خودِ قربانی \en{payload} را می‌چسباند، و تزریق‌های
پنهان را که رمزگذاری‌شده‌اند یا در وجه دیگری حمل می‌شوند از هم جدا می‌کند.
سطحِ بارگذاری \en{QR} در کار ما (بخش~\ref{sec:res:indirect}) نمونه‌ای از
ردهٔ آخر است: \en{payload} هرگز تایپ نمی‌شود و رمزگشای تصویرِ خودِ برنامه
سازوکار تحویل است. دوم، رده‌بندیِ \emph{تهدید}
آن‌ها~\cite[\en{Sec.~3.2}]{greshake2023}، شامل گردآوری اطلاعات، نفوذ و
کنترل از راه دور، محتوای دستکاری‌شده، کلاهبرداری، بدافزار و
دسترس‌پذیری، آنچه یک تزریق \emph{هست} را از آنچه \emph{به دست می‌آورد}
جدا می‌کند؛ همان تمایزی که نتایج ما وقتی موفقیتِ سطحِ متن به اجرای سمت
سرور نمی‌رسد به آن نیاز دارند (بخش~\ref{sec:disc:layers}). سوم، این
مشاهده که \en{retrieval} مسیری می‌گشاید که اغلب هیچ پالایشِ ورودی روی آن
اعمال نمی‌شود~\cite[\en{Sec.~3}]{greshake2023}، همان چیزی است که
\en{guardrail ladder} در بخش~\ref{sec:res:ladder} با ثابت نگه‌داشتنِ
پیکرهٔ حمله و جابه‌جاکردنِ \en{filter} به یک اندازه‌گیری تبدیلش می‌کند.

دو تفاوت گفتنی است. \en{Greshake} و همکاران \emph{امکان‌پذیری} را روی
چند سامانهٔ واقعی و مصنوعی نشان می‌دهند، از جمله یک گفت‌وگوی جست‌وجوی
مبتنی بر \en{GPT-4} و موتورهای تکمیل کد، و شواهدشان نمایشی است: یک حمله
یا کار می‌کرد یا نه. ما یک برنامه و یک مدل را ثابت می‌گیریم و
\emph{نرخ‌ها} را با مخرجِ اعلام‌شده گزارش می‌کنیم، که معاملهٔ یک ارزیابی
برای قابلیت انتساب است. و مصالحه در کار آن‌ها با بازرسیِ آنچه مدل کرده
قضاوت می‌شود، حال آنکه هدفِ اینجا ثابت‌هایی عرضه می‌کند --- رازِ هر سطح،
\en{PII} با شکل‌های معلوم، یک \en{routing flag} حدس‌ناپذیر برای هر کاربر
--- که همان قضاوت را در برابر \en{ground truth} قابل تصمیم می‌سازند
(بخش~\ref{sec:arch:detectors}). مقالهٔ آن‌ها همچنین یادآور می‌شود که در
زمان نگارش، \en{mitigation}های مؤثر در دسترس نبوده‌اند؛ نتایجِ
لایه‌تفکیک‌شدهٔ بخش~\ref{sec:ablation} سهمی کوچک و مشخص در همان شکاف
است، از این جهت که می‌گوید یک تکنیک کدام لایه را شکست می‌دهد، نه اینکه
یک دفاع چقدر قوی است.

\subsection{\en{Data Poisoning} و مقیاس مدل}
\label{sec:rel:poisoning}

\en{Bowen} و همکاران~\cite{bowen2025scaling} پیشینهٔ نیمهٔ مسموم‌سازی‌اند و
تفسیر آن را مقید می‌کنند. آن‌ها سه مجموعه‌دادهٔ مسموم می‌سازند، یکی برای هر
مدلِ تهدید --- \en{fine-tuning} بدخواهانه، گزینشِ ناقصِ داده، و آلوده‌سازیِ
عامدانهٔ داده~\cite[\en{Sec.~3--4}]{bowen2025scaling} --- که هر کدام یک
پیکرهٔ بی‌خطر با \num{5000} نمونه‌اند آمیخته با کسر سمی حداکثر
\num{$2\,\%$}، و نتیجه را با یک \emph{\en{learned vulnerability score}}
امتیاز می‌دهند: تغییر در یک سنجهٔ درجه‌بندی‌شدهٔ زیان‌باری، سوگیری یا
\en{backdoor} نسبت به همان مدل پیش از
\en{fine-tuning}~\cite[\en{Sec.~5.3}]{bowen2025scaling}.

دو یافتهٔ آن‌ها بر یافته‌های ما اثر می‌گذارد. مدل‌های مرزی حتی پشت یک
سامانهٔ \en{moderation} در همان کسرهای کوچک آسیب‌پذیرند: \en{GPT-4} و
\en{GPT-4o} پس از پنج \en{epoch} به امتیازهای نزدیک‌بیشینهٔ سوگیری احساسی
رسیدند، \en{GPT-4o mini} از مجموعهٔ \en{backdoor} کد در نرخ
\num{$2\,\%$} رفتار \en{sleeper agent} قابل اندازه‌گیری پیدا کرد، و
تغییر جزئیِ دادهٔ مسموم برای عبور از بررسی \en{moderation} کافی
بود~\cite[\en{Sec.~6, 6.1}]{bowen2025scaling}. افزون بر آن، روی
\num{24} مدلِ وزن‌باز از هشت خانواده در بازهٔ \num{1.5} تا \num{72}
میلیارد پارامتر، رگرسیونِ \en{learned vulnerability score} بر لگاریتم
شمار پارامترها برای مجموعه‌های \en{harmful-QA} و \en{sentiment-steering}
مثبت و از نظر آماری معنادار است~\cite[\en{Table~2}]{bowen2025scaling}:
مدل‌های بزرگ‌تر رفتار مسموم را سریع‌تر می‌آموزند، با یک خانواده
(\en{Gemma-2}) که روند وارونه نشان
می‌دهد~\cite[\en{Table~3}]{bowen2025scaling}.

سطوح مسموم‌سازیِ ما در نقطهٔ دیگری از خط لوله می‌نشینند. \en{Bowen} و
همکاران یک پیکرهٔ \en{fine-tuning} را مسموم می‌کنند و یک امتیاز
درجه‌بندی‌شده را روی یک مدل مولد می‌سنجند؛ ما آنچه را یک برنامهٔ وب متعارف
می‌بلعد مسموم می‌کنیم، یعنی یک فرم نظر که یک گام بازآموزی را تغذیه می‌کند
و سندی که به پیکرهٔ \en{retrieval} بارگذاری می‌شود، و طبقه‌بندِ زیر حمله
یک مدل قطعیِ \en{scikit-learn} است، بنابراین پاسخِ دزِ ما به‌جای آماری
دقیق است و به‌جای نرخ، به‌صورت شمارِ سطرهای تأمین‌شده به‌دست مهاجم بیان
می‌شود (بخش~\ref{sec:res:poisoning}). دو نتیجه بر سرِ این نکتهٔ عملیاتی
همداستان‌اند که مقدار بسیار کمی دادهٔ آموزشیِ در اختیار مهاجم کافی است.
جایی که کار آن‌ها برای ما تعیین‌کننده است یافتهٔ مقیاس است، و به زیان ما
عمل می‌کند: این مطالعه یک مدل یک‌میلیاردی را ثابت نگه می‌دارد، پس هر نرخ
مطلقی اینجا کرانی پایینی است که در انتهای کوچکِ روندی گرفته شده که بنا بر
شواهد آن‌ها در جهت نامطلوب برای مدافع حرکت می‌کند.
بخش~\ref{sec:disc:limitations} این پیامد را صریح بیان می‌کند به‌جای آنکه
ضمنی بگذاردش.

\subsection{آنچه این مطالعه می‌افزاید}
\label{sec:rel:gap}

میان این دو، شکافی هست که این مقاله در آن می‌نشیند. \en{Greshake} و
همکاران نشان می‌دهند کانال‌های دادهٔ یک برنامه مدل را مصالحه می‌کنند؛
\en{Bowen} و همکاران نشان می‌دهند دادهٔ آموزشیِ یک برنامه مدل را مصالحه
می‌کند. هیچ‌کدام نمی‌پرسند تیمی که \emph{مالک} چنین برنامه‌ای است باید چه
چیزی را دربارهٔ آن، در برابر سیاست خودش، و در قالبی که مهندس دیگری بتواند
بازپخش و ممیزی کند، اندازه بگیرد. پرسش اینجا همان است، و پاسخی که پیشنهاد
می‌کند روش‌شناختی است نه پدیدارشناختی: کل ارزیابی را به‌شکل \en{plugin}
برای پویشگری بیان کن که سازمان همین حالا اجرایش می‌کند؛ \en{ground truth}
هر تلاش را روی پاسخ حمل کن تا امتیازدهی دقیق باشد؛ هرگاه
\en{ground truth} غایب است به‌جای صفردادن امتناع کن؛ و هر سیگنال مقایسه
--- یک \en{detector} عمومی استاندارد، پرچمِ نشتِ خودِ برنامه، یک داورِ
اختیاری --- را به‌جای مصرف‌کردن، در کنار حکم حمل کن، تا بتوان ابزارها را
از همان اجرایی توصیف کرد که نتایج را تولید می‌کند
(بخش~\ref{sec:ablation}).
